% sec-01: Document 01, SB 1188, the California site authorization act.
\docpart{SB 1188 Site Authorization}{SB 1188 - California PDAC LLM Oncology
Clinical Trial Site Authorization and Site Establishment Act of 2026}

\section{Legislative Findings and Declarations}

The Legislature finds and declares all of the following.

\draftnote{Write ten lettered findings. Each one must carry a citation to an
author work rather than an assertion. Draw the evidence from
\mvfile{funding/pdac-funding-applications/final-apply/sections/sec-05-trial-evidence.tex}
for the simulation-to-observation chronology, from
\mvfile{funding/move-in/inputs/READMES/README-LLM-Pancreatic-Oncology-Clinical-Trial-System-Large-Documents-Funding-and-AI-Peer-Review.md}
for the twenty-paper readiness record, and from
\mvfile{funding/move-in/inputs/READMES/README-Physical-AI-Oncology-Clinical-Trial-Site-Complete-Documentation-Package.md}
for the predecessor site's operating record. Finding (a) states the disease
burden and the resectable-setting gap. Finding (j) states that California has an
opportunity to authorize the first such site and that the authorization is
conditional and revocable.}

\begin{enumerate}
\item Pancreatic ductal adenocarcinoma remains among the most lethal solid
  tumors, and the resectable setting is served by fewer trials than the
  metastatic setting. \emph{Stage 2 completes this finding with its citation.}
\item \emph{Nine further findings, to be written in full-move-in.}
\end{enumerate}

\section{Definitions}

\draftnote{Define, in lettered subdivisions: LLM-directed clinical workflow;
advisory-only output; robotic procedure suite; site safety officer; qualified
site; department; sponsor-investigator; audit trail; and verification package.
Take the advisory-only definition from
\mvfile{regulatory/Adaption-21-CFR-Part-312/} and the audit-trail definition
from \mvfile{funding/move-in/draft-move-in/sections/sec-03-sb-964-data-protection.tex},
so the two bills cannot drift.}

\section{Site Authorization}

\draftnote{Write the authorization procedure as a numbered sequence with named
timelines: application content, completeness review at thirty days, substantive
review at ninety days, conditional authorization, and the first-participant
hold. Cite \mvfile{regulatory/adaption-ich-e6r3/} for the good clinical practice
conditions the department may attach.}

\begin{mvtable}
\begin{tabularx}{\textwidth}{>{\raggedright\arraybackslash}p{1.1cm} >{\raggedright\arraybackslash}p{4.0cm} >{\raggedright\arraybackslash}p{2.2cm} >{\raggedright\arraybackslash}X}
\headrow \thead{Step} & \thead{Action} & \thead{Statutory clock} & \thead{Consequence of the department's silence}\\
\midrule
1 & Application filed & to be fixed & to be fixed \\
2 & Completeness review & to be fixed & to be fixed \\
\end{tabularx}
\end{mvtable}
\tabcapl{tab:sb1188auth}{The authorization sequence. Stage 2 fixes every clock
and states what happens when the department does not act, because a statute that
is silent on inaction is a statute with no schedule.}

\section{Operating Conditions}

\draftnote{Six conditions, each mapped to the document in this package that
implements it: staffing to document 14, premises to document 09, model
governance to document 03, safety reporting to document 13, emergency response
to document 11, and record retention to document 03.}

\section{Inspection, Enforcement, and Suspension}

\draftnote{Three deficiency classes with the department's remedy for each,
matching the class definitions written in document 06 so a site cannot be
class A under one instrument and class B under another.}

\section{Reporting and Sunset}

\draftnote{An annual report to the Legislature with named contents, and a sunset
five years from operative date unless extended. State the sunset in the same
subdivision as the report, so the report is the evidence the extension turns on.}
